\section{개요}

한국산업은행은 대한민국의 대표적인 정책금융기관으로서 산업 발전과 금융시장 안정이라는 공적 목적을 수행하기 위해 설립된 국책은행이다. :contentReference[oaicite:0]{index=0}  
일반적인 상업은행이 수익성 중심의 금융중개 기능을 수행하는 것과 달리, 한국산업은행은 국가 경제정책과 긴밀히 연계된 금융 기능을 수행한다는 점에서 본질적으로 다른 성격을 가진다.

한국산업은행의 설립 목적은 산업의 개발 및 육성, 사회기반시설 확충, 지역개발, 금융시장 안정, 그리고 지속가능한 성장 촉진에 필요한 자금을 공급·관리하는 데 있다. :contentReference[oaicite:1]{index=1}  
즉, 이는 단순한 금융기관이 아니라 국가 경제 구조를 형성하고 조정하는 정책 수단으로 이해할 필요가 있다.

이러한 특성으로 인해 한국산업은행은 투자은행(Investment Bank, IB) 기능과 정책금융 기능을 동시에 수행하는 복합적 기관으로 자리 잡았다. 대규모 기업금융, 프로젝트 파이낸싱, 인수합병(M\&A), 구조조정 지원 등은 민간 금융기관이 감당하기 어려운 영역으로, 산업은행의 핵심 역할에 해당한다. :contentReference[oaicite:2]{index=2}  

특히 경제위기 상황에서 산업은행의 역할은 더욱 두드러진다. 외환위기 이후 대기업 구조조정, 금융시장 안정화, 유동성 공급 등의 기능을 수행하며 국가 경제의 최후의 안전판 역할을 수행해왔다. 이는 산업은행이 단순한 금융기관이 아니라 거시경제 안정화 장치로 기능하고 있음을 보여준다.

그러나 이러한 공공적 기능은 동시에 구조적 긴장을 수반한다. 정책금융기관으로서 공공성과 시장 효율성 사이의 균형 문제, 정부 영향력에 따른 정치적 개입 가능성, 민간 금융기관과의 경쟁 관계 등은 산업은행을 둘러싼 지속적인 논쟁의 핵심 쟁점이다.

따라서 한국산업은행은 다음과 같은 이중적 성격을 가진다:
\begin{itemize}
    \item 공공기관으로서 국가 정책 수행 수단
    \item 금융기관으로서 시장 내 경쟁 주체
\end{itemize}

이러한 이중성은 산업은행의 기능과 역할을 이해하는 데 핵심적인 분석 틀을 제공한다.

본 문서는 한국산업은행을 단순한 기관 소개 수준을 넘어, 정책금융기관으로서의 구조적 역할과 경제적 의미를 중심으로 분석한다. 이를 위해 설립 목적과 기능, 역사적 전개, 조직 구조, 주요 사업, 정책금융 역할, 그리고 주요 논란을 체계적으로 검토한다.